%% =================================================================
%% Beber et al. IEEE Robotics and Automation Letters, Feb 2026.
%% Reconstruction of page 6 (printed p. 6).
%% Contents: Fig. 3 placeholder (three-panel trajectory comparison
%% HEDAC vs BO-EI vs BO-EIS at snapshots t = 14/28/42/56/70 s);
%% continuation of Section IV-B (robustness analysis, ergodic-metric
%% evolution); Section IV-C Comparison with BO Baselines;
%% start of Section IV-D Segmentation.
%% =================================================================

\documentclass[11pt]{article}

\usepackage[margin=0.85in]{geometry}
\usepackage{amsmath, amssymb}
\usepackage{mathrsfs}
\usepackage{xcolor}
\usepackage{fancyhdr}

\pagestyle{fancy}
\fancyhf{}
\lhead{\small 6\ \ IEEE RAL, Accepted Feb 2026}
\rhead{\small }
\cfoot{\small 6 $\to$ \arabic{page}}
\renewcommand{\headrulewidth}{0.4pt}

\setcounter{page}{1}
\setcounter{equation}{5}
\setlength{\parskip}{6pt}
\setlength{\parindent}{0pt}

\begin{document}

% ---------- Fig. 3 placeholder ----------
\begin{center}
\fbox{\begin{minipage}[c][8cm][c]{0.95\textwidth}
\centering
\textbf{Fig. 3.}\ End-effector tip trajectory during exploration of an
elasticity distribution with three stiff regions: planner comparison.\\[4pt]
\textcolor{gray}{\footnotesize
Five time snapshots in columns: $t = 14$ s, 28 s, 42 s, 56 s, 70 s.\\
Two rows: (top) Expected Information Density (EID) map; (bottom)
Estimated stiffness map.\\[3pt]
\textbf{(a)} Ergodic exploration with HEDAC (ours). The trajectory
smoothly covers the domain in proportion to EID; the cyan circle
marks trajectory start and the magenta square marks trajectory end.\\[2pt]
\textbf{(b)} BO-EI. Orange markers indicate BO-selected sampling
locations; red markers indicate points sampled along the path
between them.\\[2pt]
\textbf{(c)} BO-EIS. Same marker convention as (b); continuous
sampling augments BO's discrete selections.}
\end{minipage}}
\end{center}

\vspace{0.4em}

\textbf{Fig. 3.}\ End-effector tip trajectory during the exploration
of an elasticity distribution with three stiff regions: comparison of
different planners. The cyan circle indicates the beginning of the
trajectory, while the magenta square denotes the end of the
trajectory. In plots (b) and (c), the orange markers indicate the
locations selected by BO, whereas the red markers represent the
sampled points collected along the executed trajectory.

\vspace{0.8em}

% ---------- Section IV-B continuation ----------
time across the three scenarios, regardless of the presence of
measurement uncertainty. These results demonstrate the robustness of
the method with respect to both the number of stiff regions to be
detected and the influence of measurement noise. The RMSE, instead,
grows under both experimental conditions as the number of stiff
regions increases. This behaviour may be attributed to uncertainty
at the boundaries of the stiff regions, which are typically
under-sampled, and will be examined in greater detail in Sec.~IV-D.
As expected, we observe an increased RMSE in the second experimental
condition, attributable to the presence of the noise.

Figure~3(a) shows an example of the evolution of the agent trajectory
over time, when exploring the map with 3 stiff regions. All stiffer
regions are successfully detected within 40 s. Over time, the
ergodic metric $\alpha$ decreases from 0.79 at $t = 14$ s to 0.38 at
$t = 70$ s, reflecting improved coverage of the target distribution.
In parallel, the RMSE shows a substantial reduction, dropping from
13.75 kPa to 2.85 kPa. Intermediate values reinforce this trend: at
$t = 28$ s, $\alpha = 0.54$ and RMSE $= 11.19$ kPa; at $t = 42$ s,
$\alpha = 0.51$ and RMSE $= 3.7$ kPa; and at $t = 56$ s,
$\alpha = 0.43$ and RMSE $= 2.56$ kPa. This consistent evolution
confirms that as the exploration becomes more ergodic, the accuracy
of the reconstructed stiffness map improves accordingly.

% ---------- Section IV-C ----------
\subsection*{C.\ Comparison with BO Baselines}

In this section, we focus on the comparison of the results of our
ergodic approach with the two variants of BO described in Sec.~IV-A.
We consider the distribution with 3 stiff regions, which has proven
to be the most challenging for the ergodic exploration (see Fig.~3).
The simulation time is the sum of the simulated agent motion with a
velocity of 1 cm/s and the computation time required to determine
the next sampling location. Each simulation is conducted under two
conditions: one assuming noise-free elasticity measurements, and
another incorporating additive white noise
$\mathcal{N}(0, 6.25\;\mathrm{kPa}^{2})$, similar to Sec.~IV-B. Fixed
lenght of 400, 600, and 800 mm, were used to evaluate the performance
evolution over time, ensuring a fair comparison independently of the
method-specific stopping criteria. The average results of 100
simulations are listed in Tab.~II.

This comparison highlights that BO-EI consistently underperforms
compared to the other two methods, yielding the lowest DR and the
highest RMSE across all tested conditions. Both BO-EIS and the
ergodic strategy achieve comparable detection rates once sufficient
exploration has been performed, indicating successful identification
of the number and location of stiff regions. However, the ergodic
approach consistently provides lower reconstruction errors,
resulting in more accurate stiffness maps. At shorter travelled
distances (500 mm--650 mm), BO-EIS achieves higher detection rates,
reflecting faster initial identification of stiff regions. This
behaviour is expected given the exploitative nature of BO-based
strategies, which prioritise sampling high-uncertainty locations.
However, in clinical settings, once suspicious regions are
identified, accurate post-detection refinement becomes critical.
Therefore, it is more relevant to compare the approaches once
detection performance is maximised (800 mm, DR $\approx 99$--100\%),
and differences between the methods are primarily reflected in
reconstruction accuracy and execution duration. In this regime, the
ergodic strategy achieves lower RMSE while requiring substantially
less execution time than BO-EIS under the same travelled-distance
budget, both in noise-free and noisy conditions. These results
indicate that, in the post-detection phase, ergodic exploration
enables more efficient refinement of stiffness maps, yielding
improved accuracy per unit motion and reduced execution time under
equal motion cost. Compared to BO-based approaches, the ergodic
strategy also exhibits greater robustness to noise, maintaining a
consistent performance advantage across all evaluated conditions.
Overall, these findings confirm the benefit of ergodic exploration
in balancing detection reliability, reconstruction accuracy, and
execution efficiency once suspicious regions have been identified.

% ---------- Section IV-D ----------
\subsection*{D.\ Segmentation}

The spatial accuracy of the proposed method was evaluated based on
its ability to segment stiff regions within the reconstructed
stiffness maps and compared against the BO. To this end, we designed
two simulated stiffness distributions: the first one with two
inclusions, one irregularly shaped (Cluster 1) and one circular
(Cluster 2), while the second distribution consisted of a single
donut-shaped inclusion (Cluster 3), as

\end{document}
